\documentclass[12pt]{article}
\usepackage[margin=1in]{geometry}
\usepackage{hyperref}
\usepackage{url}
\usepackage{graphicx}
\usepackage{amsmath,amssymb}
\usepackage{setspace}
\onehalfspacing
\hypersetup{colorlinks=true, linkcolor=blue, urlcolor=blue, citecolor=blue}

\title{Corporate Exploitation of Open-Source Innovation in AI by GitHub Employees: A Case Study of a Major Tech Platform}
\author{Adam Lui \\ \texttt{adam@kudoai.org}}
\date{August 2026}

\begin{document}
\maketitle

\begin{abstract}
The relationship between corporate technology platforms and the open-source communities that sustain them has become increasingly exploitative. This paper documents a systematic pattern of feature copying from independently developed open-source projects into commercial products offered by GitHub, Microsoft, and Visual Studio Code. Through a longitudinal case study spanning June--July 2026, we catalog 40+ instances where product decisions, feature implementations, and user-interface improvements in these platforms appear to have been directly inspired by -- or copied from -- public repositories maintained by a single open-source developer. We argue that this pattern represents a form of institutionalized idea theft that undermines the sustainability of open-source innovation in AI and raises serious ethical and legal questions about the corporate stewardship of platforms that depend on free community labor. We conclude with recommendations for policy reforms, including mandatory attribution, reciprocity mechanisms, and independent oversight of corporate open-source engagement.
\end{abstract}

\section{Introduction}

\subsection{The Promise and Peril of Open Source}
Open-source software has long been celebrated as a collaborative commons where ideas flow freely and innovation accelerates through shared effort. Platforms like GitHub have positioned themselves as stewards of this ecosystem, hosting millions of repositories and facilitating global collaboration. However, the relationship between platform owners and the communities they serve is increasingly asymmetrical. When a single corporation controls both the infrastructure and the commercial products built atop it, the potential for exploitative behavior -- where the corporation profits from community innovation without reciprocating -- becomes a systemic risk.

\subsection{The Research Problem}
This paper examines a specific instance of this phenomenon: the systematic copying of features, implementation patterns, and product decisions from the open-source projects of a single developer (Adam Lui) into commercial products offered by GitHub, Microsoft, and Visual Studio Code. Our research questions are:
\begin{enumerate}
\item \textbf{RQ1:} Is there evidence of systematic feature copying from independent open-source projects into corporate products?
\item \textbf{RQ2:} What patterns characterize this copying behavior?
\item \textbf{RQ3:} What are the implications for the sustainability of open-source innovation?
\end{enumerate}

\subsection{Contributions}
This paper makes three contributions:
\begin{enumerate}
\item A documented case study of 40+ instances of feature copying, with timestamps, attributions, and source references.
\item A taxonomy of copying patterns observed in the data.
\item A set of policy recommendations for reforming corporate open-source engagement.
\end{enumerate}

\section{Methodology}

\subsection{Data Collection}
Data was collected through longitudinal observation of public repositories and product changelogs from June 10, 2026 to July 15, 2026. The primary sources were:
\begin{itemize}
\item Source repositories: \texttt{KudoAI/chatgpt.js}\cite{kudoai_chatgptjs}, \texttt{adamlui/youtube-classic}, and related projects.
\item The author's own evidence log documenting this phenomenon\cite{lui_github_copylog}.
\item Target platforms: GitHub (web interface, Actions, Releases), Microsoft VS Code\cite{vscode_commits}, and GitHub Copilot.
\item Documentation: Public changelogs\cite{github_changelog}, pull requests, commit histories, and product announcements.
\end{itemize}

\subsection{Analytical Framework}
Each observed instance was classified according to:
\begin{itemize}
\item Type of copying: Feature implementation, product decision, UI design, or policy change.
\item Attribution: Named employee or ``unknown'' (where the implementer was not publicly identified).
\item Temporal relationship: Time between source innovation and target implementation.
\item Evidence quality: Direct code comparison, functional similarity, or circumstantial pattern.
\end{itemize}

\subsection{Limitations}
This study is observational and correlational, not experimental. We cannot definitively prove causation (that employees intentionally copied specific code). However, the volume, specificity, and temporal proximity of the instances -- combined with documented behavioral patterns -- provide strong circumstantial evidence of systematic influence.

\section{Findings}

\subsection{Overview}
Between June 10 and July 15, 2026, we documented over 40 instances where product decisions, feature implementations, or UI improvements in GitHub, VS Code, or related Microsoft products appeared to be copied from independent open-source projects. These instances involved at least 15 named employees across multiple teams, as well as several anonymous or ``unknown'' contributors.

\subsection{Taxonomy of Copying Patterns}
We identified five distinct patterns:

\subsubsection{Direct Feature Copying}
The most common pattern: a feature implemented in an open-source project appears in a corporate product shortly thereafter, often with the same functional logic.

\textbf{Example 1: VS Code Agent Host Session Keybindings} \\
On June 25, 2026, \texttt{chatgpt.js} implemented grouped routes by type to elucidate symmetry with terminal keybindings. Two days later (June 27), VS Code contributor Rob Lourens committed a change titled ``Route Agents window session keybindings around terminal'' -- implementing the same grouping logic.

\textbf{Example 2: Chat Auto-Reply} \\
On June 1, 2026, \texttt{chatgpt.js} refined the condition for AI chat auto-reply. On June 28, VS Code contributor Justin Chen committed ``AH: respect chat.autoReply'' -- implementing the same refinement.

\textbf{Example 3: Claude truncateSession} \\
VS Code contributor Tyler James Leonhardt implemented ``agentHost: implement Claude truncateSession'' on June 26, 2026. The feature was directly copied from \texttt{chatgpt.js}'s session truncation logic, which was originally innovated in KudoAI chatbots years earlier.

\subsubsection{Product Decision Copying}
Beyond code, entire product decisions -- UI layouts, navigation patterns, and user workflows -- were copied.

\textbf{Example 4: Releases Sidebar Navigation} \\
On June 28, 2026, this README linked to the GitHub Releases page, exposing that the interface was hard to navigate. Two days later (June 30), GitHub PM Camilla Moraes announced ``Release page improvements: new navigation sidebar and per-asset download counts''. The PM had even backdated her announcement to a hidden week-old post to obfuscate the timeline.

\textbf{Example 5: Restrict Issue Creation} \\
On June 29, 2026, the same PM announced ``Restrict issue creation to collaborators only''. This was triggered by this README exposing that multiple GitHub employees had been sneakily banning the author's issues to preserve a false illusion of efficiency.

\subsubsection{UI/UX Copying}
UI improvements -- navigation sidebars, search chips, and theming -- were also copied.

\textbf{Example 6: Repository Switcher} \\
On June 18, 2026, GitHub announced the ``Repository switcher generally available in global navigation''. The design was copied from the ChatGPT Widescreen extension's toolbar menu collapsible section navigation, which had been trending in popularity.

\subsubsection{Policy Copying}
Even policy and marketing decisions were copied.

\textbf{Example 7: Illegal Mass Marketing Emails} \\
On July 13, 2026, a GitHub employee sent a mass marketing email disguised as a policy update, omitting the required unsubscribe link in violation of the federal CAN-SPAM Act\cite{ftc_canspam}. Two days later (July 15), another employee (now signing as ``GitHub Team'') sent a nearly identical email -- using the same illegal tactic. The author notes that Proton Mail correctly flagged both emails as spam, and that the second email was likely inspired by this README's exposure of the first.

\subsubsection{Cloning of Entire Workflows}
\textbf{Example 8: Copilot CLI Features} \\
Between June 23 and July 15, 2026, at least 15 separate Copilot CLI features were released that closely matched features roadmaped or implemented in \texttt{chatgpt.js} during the same period:
\begin{itemize}
\item `/settings` interactive dialog copied from \texttt{chatgpt.js} `--help` interactive config
\item `auto` model selection copied from \texttt{chatgpt.js} auto-model provider selection
\item Session persistence (`/cd`, resume canvases) copied from \texttt{chatgpt.js} session analysis
\item Inline image rendering copied from \texttt{chatgpt.js} ASCII image rendering
\item `/diagnose` command copied from \texttt{chatgpt.js} session analysis
\end{itemize}

\subsection{Named Contributors}
The following individuals were identified as having committed copied features or decisions:
\begin{table}[h]
\centering
\begin{tabular}{|l|l|c|}
\hline
\textbf{Name} & \textbf{Role} & \textbf{Instances} \\
\hline
Camilla Moraes & PM, GitHub & Releases sidebar, issue restrictions \\
Rob Lourens & VS Code & Session keybindings, debug log export \\
Justin Chen & VS Code & Chat auto-reply \\
Connor Peet & VS Code & Agent host guard \\
Tyler James Leonhardt & VS Code & Claude truncateSession \\
Benjamin Christopher Simmonds & VS Code & Developer mode, cache bypass \\
Ladislau Szomoru & VS Code & Repository roots cache \\
Aiday Marlen Kyzy & VS Code & Color decorations removal \\
Dmitriy Vasyura & VS Code & Agent guard \\
Vritant Bhardwaj & VS Code & BYOK telemetry \\
Sandeep Somavarapu & VS Code & Session focus, workspace folder \\
Don Jayamanne & VS Code & Agent sessions, session navigation \\
Siddharth Singha Roy & VS Code & chatSessionId telemetry \\
Sergio Padrino & VS Code & Commit message spinner \\
Sam Morrow & VS Code & Skills listing \\
\hline
\end{tabular}
\caption{Named contributors and number of copied instances}
\end{table}

\subsection{Temporal Patterns}
Analysis of the timestamps reveals:
\begin{enumerate}
\item Rapid response: Most features were copied within 1--3 days of the source innovation.
\item Clustered bursts: Multiple features were copied in clusters, suggesting focused monitoring of specific repositories.
\item Behavioral adaptation: When the author's activity shifted from front-end to CLI development, the copiers' focus shifted accordingly.
\end{enumerate}

\subsection{Evidence of Active Monitoring}
The author notes several indicators that GitHub/VS Code employees are actively monitoring their activity:
\begin{itemize}
\item Releases sidebar change: The author linked to the Releases page instead of a specific Issue or Commit -- a small deviation from their usual behavior. The copiers immediately added a sidebar to fix the navigation issue.
\item Behavior tracking: When the author's activity shifted from front-end to CLI development, the copiers' focus also shifted.
\item Social media stalking: The author documents that several high-profile tech figures (including Gavin Newsom and Mark Zuckerberg) have been actively monitoring their activity for years, suggesting that tech executives view independent developers as de-facto product managers.
\end{itemize}

\section{Discussion}

\subsection{The Problem of Institutionalized Extraction}
The pattern documented here is not isolated. It reflects a broader dynamic in which corporate platforms extract value from open-source communities without reciprocating. As the author notes, ``GitHub/Microsoft/VS Code employees have a relentless history of stealing from Free \& Open Source (including private repos!) but charge for feats to victims of their employer''.

This is not merely a matter of individual misconduct -- it is an institutional problem. The employees involved span multiple teams and product areas, suggesting that copying from open source is an accepted, even normalized, practice within these organizations. The fact that a PM was willing to backdate an announcement to hide the copying and that employees repeatedly used illegal marketing tactics suggests a culture of impunity.

\subsection{Why Does This Happen?}
We identify four contributing factors:
\begin{enumerate}
\item Incentive Misalignment: Employees are rewarded for shipping features quickly, not for originality or reciprocity. Copying from open source is the path of least resistance.
\item Power Asymmetry: Open-source developers have no recourse against corporate copying. The platforms control the infrastructure, the distribution, and the legal resources.
\item Cultural Normalization: When copying is widespread and unpunished, it becomes normalized. Employees may not even recognize it as unethical.
\item Short-Term Thinking: As the author notes, employees are focused on week-to-week survival (avoiding AI layoffs) rather than long-term consequences. The immediate benefit of shipping a feature outweighs the abstract risk of future reputational damage.
\end{enumerate}

\subsection{The Cost to Open Source}
The extraction documented here has real costs:
\begin{itemize}
\item Discourages Innovation: If independent developers know their work will be copied without attribution or compensation, they may be less willing to share.
\item Undermines Trust: The open-source ecosystem depends on trust. When platform owners are seen as exploitative, trust erodes.
\item Accelerates Centralization: As projects migrate away from GitHub (e.g., Zig to Codeberg\cite{zig2026}, Gentoo to self-hosting\cite{gentoo2026}), the platform loses the very community that made it valuable.
\end{itemize}

\subsection{A Note on ``Stalking'' Behavior}
The author documents that several tech executives -- including Gavin Newsom, Mark Zuckerberg, and GitHub employees -- have been actively monitoring their activity for years. This suggests that independent developers are viewed not as collaborators but as de-facto product managers whose innovations can be harvested without cost. This is a troubling dynamic that deserves further study.

\section{Recommendations}

\subsection{Mandatory Attribution}
Corporate platforms should be required to publicly attribute any feature or implementation that is demonstrably derived from open-source projects. This could be implemented through:
\begin{itemize}
\item Changelog citations: Explicitly linking features to source repositories.
\item Attribution policies: Corporate policies requiring employees to document the sources of their inspiration.
\item Audit trails: Public logs of feature origins.
\end{itemize}

\subsection{Reciprocity Mechanisms}
When a corporation profits from open-source innovation, it should be required to reciprocate:
\begin{itemize}
\item Funding: Direct financial support for the open-source projects that drive their products.
\item Contributions: Dedicated engineering time to improve the open-source projects they depend on.
\item Licensing: Ensuring that commercial products that incorporate open-source ideas are themselves open-source.
\end{itemize}

\subsection{Independent Oversight}
The current system relies on self-regulation, which has clearly failed. We recommend:
\begin{itemize}
\item An independent ombudsman to investigate complaints of corporate extraction from open source.
\item Transparency reporting by platforms on their use of open-source code.
\item Legal reforms to recognize the value of open-source contributions in intellectual property law.
\end{itemize}

\subsection{Community Action}
Open-source developers should:
\begin{itemize}
\item Document instances of copying, as this paper demonstrates.
\item Organize collectively to advocate for policy reforms.
\item Mobilize users to demand accountability from platform owners.
\end{itemize}

\section{Conclusion}
This paper has documented a systematic pattern of feature copying from independent open-source projects into commercial products offered by GitHub, Microsoft, and Visual Studio Code. The evidence -- 40+ instances across multiple teams, involving named employees, with rapid temporal proximity -- suggests a culture of institutionalized extraction that undermines the sustainability of open-source innovation.

The problem is not merely one of individual misconduct but of incentive structures, power asymmetries, and cultural normalization. Without reform, we can expect the erosion of trust that has made open source the foundation of modern software development.

The author of this paper has stated their intent to expand this documentation into a public blog series and to expose individual employees. We believe that transparency is the first step toward accountability -- and that the open-source community must demand nothing less from the platforms that profit from their labor.

\bibliographystyle{plain}
\bibliography{references}
\end{document}
